% v22 research notes: retain only material useful for future work.
\subsection{三高同時非負性：低行二変数までの有限化}
\label{notes:hd22}

\paragraph{現在の到達点}

第21版の五固定データ $(s,t,x,y,z)$ ごとの $b$ 有限化は，
第22版で固定 $(s,t)$ ごとの全 $x,y,z,b$ の有限復元へ強まった．
新しい主モジュールは
\nolinkurl{high_d_fixed_low_row_finiteness_v22.tex} である．
低行 $(-s,-t,s+t-1,0)$ の二整数自体の一様上限，
および三高同時非負模型の一般不存在は未解決である．

利用する記号と核となる式だけを残す．
\[
 u=s+t-1,\ h=s-x\ge1,\ B=b-1\ge1,\
 F=B+ux+sz,\ m=F^2+F+1.
\]
固定 $(s,t,x,z,b)$ において
\[
 a=uh-B(F+1),\qquad
 n_0=(F+1)(Bt+h)+F(u+z)+htz,
 \qquad mT=n_0+ay.
\]
$K=sT-b>0$ と $y\ge0$ が
\[
 \boxed{F\le M:=st+(2u+1)h-2}
\]
を与える．上限の証明は $T$ の正整数性を過剰に使わず，entry $K$ の正値性を
直接使う．$F$ は整数なので，$F\ge M+1$ の排除から表示した上限が従う．
さらに $b-1+sz\le M-ux$ より
\[
 (3u+1)x\le st+(2u+1)s-3
\]
も従う．これは $x\le s-1$ の追加の絞り込みとして使える．

残余式を $y$ の多項式にして
\[
 \mathcal R(Y)=suaY^2+(aL-C)Y+C_0
\]
とすると，$L,C>0$ であり，$a=0$ でも一次係数が $-C<0$ となる．
従って各 $(s,t,x,z,b)$ に非負整数根は高々二個，$a\le0$ なら高々一個である．
二次式の係数と $C_0$ の計算法は公開用証明に明示した．
根は判別式の整数平方根と正確な割り算で全て求められ，$y$ の探索上限は不要である．

\paragraph{完全復元で残す検査}

二次式の根であることは元の第一残余そのものの零を意味する．
ただしそれだけでは整数模型ではない．必ず
\[
 d=x+y+z\ge1,\quad \alpha\ge k\ge b+1,\quad m\mid FE-J,
\]
三尺度の正値性，六つの符号量の正値性を検査する．
最後に元の $q,P$ を組み立て，row sums，六つの pair products，
$D=(x,y,z,-d)$，正 weights に対する十二の gaps を再計算すれば，
導出式の取り違えも検出できる．$j=t-y$ の非負性は仮定しない．

第21版の零ノルム分類は新手順の例外処理として不要になったが，
$(s,t,x,y,z)=(3,3,0,0,1)$ と $(2,2,1,1,1)$ の両族は
元残余を捨ててはいけないことを示す回帰例として残した．
前者の尺度 $(3,3,6)$ は $b$ によらず整数だが，
残余 $27-b(b+5)$ は $b\ge2$ の整数で零にならない．

\paragraph{再現可能な厳密計算と次の分岐}

標準ライブラリのみの
\nolinkurl{verify_high_d_low_row_finiteness_v22.py}
は二十四の多項式恒等式を係数ごとに確認する．
\texttt{--fixed s t} はその低行の全模型を完全復元し，
\texttt{--scan 12} は $1\le s,t\le12$ の全 $144$ 低行に同じ処理を行う．
探索箱がかかるのは $s,t$ だけであり，$D$ や分母は全範囲である．
結果は \nolinkurl{high_d_low_row_finiteness_verification_v22.json} にある．
全五組候補 $2{,}004{,}015$，$a=0$ の一次式 $219$，
非負整数残余根六個で，全て $d=0$ のため対象外だった．
零 $D$ の六根の $(s,t,b)$ は
\[
 (3,4,6),(4,3,6),(6,9,13),(9,6,13),(10,12,19),(12,10,19).
\]
これらの尺度は有理数であり，整数の零 $D$ 模型の反例ではない．

次の有力課題は，無上界の $(s,t)$ に対して残余二次式の非負整数根と
単一合同条件を同時に扱うことである．有限復元を使えば，
単純な分母・高係数 cutoff の拡大を避け，低行の構造に沿った反例探索ができる．
また観測範囲では $d\ge1$ の整数残余根が合同条件を課す前から存在しなかった．
これは，尺度の整数性を外した有理復元段階にもより強い障害がある可能性を示すが，
現時点では観測であって定理ではない．
判別式は
\[
 \mathcal D_y=(aL-C)^2-4suaC_0
 =C^2+a\bigl(aL^2-2LC-4suC_0\bigr)
\]
である．一般の全展開では扱いやすい積を得ていない．
次は $a$ の符号と $a=0$ の退化面を先に分け，
整数平方間の評価と合同条件を組み合わせる方がよい．
$a=0$ は $B(F+1)=uh$，すなわち
\[
 u(s-bx)=B(b+sz)
\]
を要求するので，$s-bx>0$ という追加制約も得られる．
この面で残余は $-CY+C_0$ であり，$y=C_0/C$ の整数性へ直ちに帰着する．

判別式とは別に，$s=t$ かつ $x=y$ の対称 sector は無上界で排除できた．
このとき $J=-E$，単一合同条件から $n=E/m\in\mathbb Z$，
$S=h+(F+1)n$，$T=h-Fn$ となる．
$F\ge h+1$ により正値性は $n=0$ を強制し，既知の零ノルム排除に矛盾する．
この証明は整数性と正値性の両方を使うので，
一般 sector でも $E+J=\ell(sy-tx)+h-j$ の小さな値から類似の格子拘束を
引き出せるかを検討する価値がある．
一般不存在を目標に固定せず，低行に沿った無限族・合同条件の反例・
残余式の因数分解が現れた場合は，その構造の記述へ研究計画を切り替える．
